\scanpage{97}
\mspage{\unsure{\ldots)|3}}
concerning this secondary structure. The first thing I could think of is the ``pentagon axiom'' for the associativity data, which occurs when looking at associativities for the compositum (for $\overset{\deleted{\unsure{i}}i}{*}_\ell$ say) of four factors. Here again\deleted{\unsure{,}} the first reflex would be to write down, as usual, an \uline{equality} for \deleted{the} two compositions of associativity isomorphisms, exhibited in the pentagon diagram. One suspects however that such equality should, again, be replaced by a ``homotopy''-arrow, which now appears as a kind of ``ternary'' structure - before even having exhausted the list of coherence \handwritten{``}relations\handwritten{''} one could think of with respect to the secondary structure ! Here one seems caught at first sight in an infinite chain of \deleted{increasingly}\inserted{\handwritten{ever}} \handwritten{``}higher\handwritten{''}, and presumably\handwritten{,} messier structures, where one is going to get hopelessly lost, unless one discovers some simple guiding principle for sheding some clarity in the mess.

\handwritten{3.}\quad I thought of writing you mainly because I believe that, if anybody, you should know if the kind of structure I am looking for has been worked out - maybe even \uline{you} did ? In this respect, I vaguely remember that you had a description of $n$-categories in terms of $n$-semisimplicial sets, satisfying certain exactness conditions, in much the same way as an ordinary category can be interpreted, via its ``nerve'', as a particular type of semisimplicial set. But I have no idea if your definition applied only for describing $n$-categories with strict associativities, or not.

Still some comments in the spirit of your axiomatics of homotopical algebra - in order to make the question I am proposing more seducing maybe to you ! One comment is that presumably, the category of $\infty$-groupoids (which is still to be defined) is a ``model category'' for the usual homotopy category; this would be at any rate one plausible way to make explicit the intuition referred to before, that a homotopy type is ``essentially the same'' as an $\infty$-groupoid up to $\infty$-equivalence. The other comment: the construction of the fundamental $\infty$-groupoid of a space, disregarding for the time being the question of working out in full the pertinent structure on this messy object, can be paraphrased in any model category in your sense, and yields a functor from this category to the category of $\infty$-groupoids, and hence (by geometric realization, or by \unsure{localisation|localization}) also to the usual homotopy category. Was this functor obvious beforehand ? It is of a non-trivial nature only when the model category is \uline{not} pointed - as a matter of fact the whole construction can be carried out canonically, in terms of a ``cylinder object'' $\unsure{I|\mathbb{I}}$ for the final object $e$ of the model category\handwritten{, playing the role of the unit segment.}
% archive sheet number: 96
